% Responsable: Joseph. Rutas, autenticación, autorización, rate limiting, endpoints y justificación.
\section{API Gateway y comunicación síncrona}
\subsection{Diseño del API Gateway}

El API Gateway será el punto de entrada del cliente web del SGRA.
Recibirá las solicitudes por HTTPS y las enviará al microservicio
correspondiente según la ruta y el método HTTP. Por ejemplo, una
solicitud de refuerzo se dirige al servicio Refuerzos, mientras que
una consulta de aulas disponibles se dirige al servicio Espacios.

Para el diagrama C4 se proponen tres componentes del Gateway:
el enrutador, el filtro de acceso y el filtro de control de tráfico.
El enrutador dirige cada solicitud. El filtro de acceso valida el
token en las rutas protegidas. El filtro de tráfico limita la
cantidad de solicitudes que recibe el Gateway. Cada microservicio
comprueba si el usuario tiene permiso sobre el recurso solicitado.

\subsection{Rutas del API Gateway}

La Tabla~\ref{tab:rutas-gateway} muestra las principales solicitudes
del cliente web y el microservicio que atiende cada una.

\small
\begin{longtable}{@{}p{.40\linewidth}p{.20\linewidth}p{.34\linewidth}@{}}
\caption{Rutas principales del API Gateway del SGRA.}
\label{tab:rutas-gateway}\\
\toprule
\textbf{Método y ruta} & \textbf{Servicio} & \textbf{Operación} \\
\midrule
\endfirsthead
\toprule
\textbf{Método y ruta} & \textbf{Servicio} & \textbf{Operación} \\
\midrule
\endhead
\bottomrule
\endlastfoot

\texttt{POST} \path|/api/auth/login| &
Identidad y acceso & Validar credenciales y emitir el token. \\

\texttt{GET} \path|/api/asignaturas| &
Académico & Consultar las asignaturas del período. \\

\texttt{GET} \path|/api/solicitudes| &
Refuerzos & Listar las solicitudes que puede ver el usuario. \\

\texttt{POST} \path|/api/solicitudes| &
Refuerzos & Registrar una solicitud de refuerzo. \\

\texttt{POST} \path|/api/sesiones| &
Refuerzos & Programar una sesión de refuerzo. \\

\texttt{GET} \path|/api/sesiones/{id}| &
Refuerzos & Consultar el detalle y el estado de una sesión. \\

\texttt{PATCH} \path|/api/sesiones/{id}| &
Refuerzos & Reprogramar una sesión. \\

\texttt{PUT} \path|/api/sesiones/{id}/asistencia| &
Refuerzos & Registrar o actualizar la lista de asistencia. \\

\texttt{GET} \path|/api/espacios/disponibilidad| &
Espacios & Consultar las aulas o laboratorios disponibles. \\

\texttt{GET} \path|/api/reportes| &
Seguimiento & Consultar los reportes de coordinación académica. \\

\texttt{PATCH} \path|/api/admin/parametros| &
Administración & Actualizar los parámetros generales del sistema. \\

\end{longtable}
\normalsize

\subsection{Llamadas síncronas entre microservicios}

\textbf{De Refuerzos a Académico.}
Antes de guardar una solicitud, Refuerzos consulta
\texttt{GET} \path|/interno/matriculas/vigentes|.
La respuesta de Académico determina si el estudiante está matriculado
en la asignatura para la que solicita el refuerzo.

\textbf{De Refuerzos a Espacios.}
Al programar o reprogramar una sesión presencial, Refuerzos usa
\texttt{PUT} \path|/interno/reservas/{idSesion}| para pedir la reserva.
Espacios comprueba la disponibilidad y responde si pudo reservar.
Si se agota el tiempo de espera, el cambio queda pendiente.
Refuerzos consulta \texttt{GET} \path|/interno/reservas/{idSesion}|
para comprobar si la reserva se completó antes de volver a solicitarla.

Las rutas \path|/interno/| son para la comunicación entre servicios.
Una vez confirmado el cambio, el aviso al usuario se procesa mediante
un evento que consume Notificaciones.

\subsection{Selección de REST y relación con los RNF}

Se propone REST sobre HTTPS con JSON para las solicitudes del cliente
web y las llamadas síncronas entre microservicios. Las operaciones del
SGRA consisten en consultar o modificar recursos como solicitudes,
sesiones y reservas. Angular puede consumir estas respuestas JSON con
su cliente HTTP, lo que facilita la integración con el navegador
prevista en el RNF04 \cite{sgra-angular-http}.

gRPC también permite realizar llamadas síncronas y definir contratos
tipados con Protocol Buffers. Para usarlo desde el navegador habría que
integrar gRPC-Web. Dado que las operaciones identificadas se resuelven
con una petición y una respuesta, se propone mantener REST en ambos
canales \cite{sgra-grpc-web}.

\begin{itemize}
    \item \textbf{RNF01, rendimiento.} Para la meta de responder en
    un máximo de 5 segundos con hasta 100 usuarios concurrentes,
    el Gateway asignará un tiempo límite a las llamadas internas.
    Esta meta se comprobará con pruebas de carga de extremo a extremo.

    \item \textbf{RNF02, disponibilidad.} Para la meta del 99,5\,\%
    mensual se proponen réplicas del Gateway detrás de un balanceador
    con comprobaciones de salud. El Gateway podrá devolver HTTP~503
    ante una indisponibilidad temporal o HTTP~504 si el servicio
    consultado supera el tiempo de espera \cite{sgra-http-semantics}.

    \item \textbf{RNF05 y RNF06, auditoría e integridad.} Cada servicio
    registrará las acciones críticas y verificará que el usuario tenga
    permiso sobre el recurso antes de modificarlo.
\end{itemize}